\documentclass[a4paper]{article}
\usepackage[pdftex]{graphicx}
\usepackage[utf8]{inputenc}
\usepackage{enumerate}
\usepackage{siunitx}
\sisetup{locale=DE}
\usepackage{href-ul}
\hypersetup{
	colorlinks=true,
	linkcolor=blue,
	urlcolor=blue} 
\usepackage{icomma}
\usepackage{amssymb}
\usepackage{geometry}
\geometry{a4paper, top=15mm, left=15mm, right=15mm, bottom=15mm,
	headsep=10mm, footskip=12mm}

\begin{document}
{\bf Newton-Mechanik, beschleunigte Bewegung}
\begin{enumerate}[1.]
{\bf \item Berechne die fehlenden Werte:}\\
Die Szenarien in der Tabelle beziehen sich auf Objekte, die aus dem Stillstand beschleunigt werden.

\renewcommand{\arraystretch}{2}
\begin{tabular}{|p{4 cm}|p{1.6 cm}|p{1.8 cm}|p{1.8 cm}|p{1.8 cm}|p{1.8 cm}|p{2.0 cm}|}
	\hline
	 &a) & b ) & c) & d) & e) & f) \\
	 & Fahrrad & Auto & Freier Fall & Rakete & Gewehrkugel & Elektron \\
	\hline
	Endgeschwindigkeit v & $\SI{10}{\meter\over \second}$& & & $\SI{8,0}{\kilo\meter\over \second}$& $\SI{350}{\meter\over \second}$& \\
	\hline
	Beschleunigungsstrecke s & & & &$\SI{500}{\kilo\meter}$ & $\SI{1,0}{\meter}$&$\SI{10}{\centi\meter}$ \\
	\hline
	Zeitintervall $\Delta t$  & $\SI{8,0}{\second}$ & & $\SI{3,0}{\second}$& & & \\
	\hline
	Beschleunigung a & & $\SI{3,5}{\meter\over{\second^2}}$ &$\SI{9,81}{\meter\over{\second^2}}$ & & & \\
	\hline
	Masse m & $\SI{105}{\kilogram}$ & & & $\SI{5000}{\kilogram}$ &&  $\SI{9,1e-31}{\kilogram}$\\
	\hline
	Trägheitskraft F & & $\SI{2975}{\newton}$ &$\SI{0,98}{\newton}$ & & & \\
	\hline
	Beschleunigungsarbeit $W$ & & $\SI{595}{\kilo\joule}$ & & &$\SI{612}{\joule}$ &$\SI{2,2e-18}{\joule}$ \\
	\hline
\end{tabular}

{\bf \item Berechne die fehlenden Werte:}\\
Die Szenarien in der Tabelle beziehen sich auf Objekte, die eine Anfangsgeschwindigkeit besitzen und beschleunigt oder abgebremst werden.

\renewcommand{\arraystretch}{2}
\begin{tabular}{|p{5 cm}|p{2 cm}|p{2 cm}|p{2 cm}|p{2 cm}|p{2 cm}|}
\hline
	& g) & h) & i) & j)&k) \\	
\hline
Anfangsgeschwindigkeit $v_1$ & $\SI{5,0}{\meter\over\second}$ & $\SI{15}{\meter\over\second}$& $\SI{10}{\meter\over\second}$& &\\
\hline
Endgeschwindigkeit $v_2$ & $\SI{14}{\meter\over\second}$ & & $\SI{0,0}{\meter\over\second}$ &$\SI{18}{\meter\over\second}$ & $ \SI{0,0}{\meter\over\second}$\\
\hline
Beschleunigungsstrecke s & & $\SI{100}{\meter}$&  & & $\SI{5,0}{\milli\meter}$\\
\hline
Zeitintervall $\Delta t$  & & & $\SI{1,0}{\second}$ & $\SI{10}{\second}$ &\\
\hline
Beschleunigung a & $\SI{9,81}{\meter\over{\second^2}}$ & &  &  & $\SI{2500}{\meter\over{\second^2}}$\\
\hline
Masse m  & $\SI{0,40}{\kilogram}$ & $\SI{900}{\kilogram}$ & &$\SI{250}{\kilogram}$ &\\
\hline
Trägheitskraft F & & $\SI{3,0}{\kilo\newton}$& & $\SI{800}{\newton}$ &\\
\hline
Beschleunigungsarbeit $\Delta W$  & & &$\SI{4,0}{\joule}$ & & $\SI{3,75}{\joule}$\\
\hline
Vorgang (bitte wählen) & & & & &\\
\hline
\end{tabular}

\begin{enumerate}[I.]
	\item Senkrechter Wurf nach oben
	\item Senkrechter Wurf nach unten
	\item Auto beschleunigt
	\item Hammer auf Nagel
	\item Motorrad bremst
\end{enumerate}

\end{enumerate} 
\fbox{
	\begin{minipage}{0.5\textwidth}
		Zur Lösung bitte \href{https://www.okuyakl.de/phys/p10nmcL083/ll083.pdf}{hier klicken} oder den QR-Code scannen.\\
	Weitere Arbeitsblätter gibt es unter 
	
	\href{https://www.okuyakl.de}{www.okuyakl.de}
	\end{minipage}
	\hfill
	\begin{minipage}{0.4\textwidth}
		\includegraphics[width=1.5 cm]{../../viecher/zwe03}
		\includegraphics[width=3 cm]{qr083}
		\includegraphics[width=2 cm]{../../viecher/afanticon1}
		
	\end{minipage}}
\end{document}